\title{DeepSeekMath: Pushing the Limits of Mathematical Reasoning in Open Language Models}

\begin{abstract}
Mathematical reasoning remains a formidable task for language models, primarily due to its intricacy and the rigorous structure involved. This work presents DeepSeekMath~7B, an extension of DeepSeek-Coder-Base-v1.5 7B that is further pre-trained with 120B mathematics-centric tokens curated from Common Crawl, complemented by both natural language and programming code. DeepSeekMath~7B secures a notable 51.7\% on the MATH benchmark at competition standard, all without dependence on auxiliary toolkits or ensemble methods—achieving a level of accuracy comparable to leading systems like Gemini-Ultra and GPT-4. By applying self-consistency with 64 samples, DeepSeekMath~7B attains a score of 60.9\% on MATH. The model’s advancements in mathematical reasoning stem from two core innovations: (1) the exploitation of publicly accessible web resources using a carefully constructed data curation framework, and (2) the development of Group Relative Policy Optimization (GRPO)—an adaptation of Proximal Policy Optimization (PPO)—which both bolsters mathematical reasoning proficiency and improves PPO’s memory efficiency.
\end{abstract}

While DeepSeekMath~delivers strong results in quantitative reasoning tasks, its abilities with geometry and theorem-proving lag behind closed-source models. For example, during preliminary tests, the model was unable to solve triangle and ellipse-related questions, suggesting a possible bias in data curation during the pre-training and fine-tuning stages. Moreover, due to model size constraints, DeepSeekMath~does not match GPT-4’s performance in few-shot scenarios: whereas GPT-4 benefits from additional example inputs, DeepSeekMath~shows little change between zero-shot and few-shot results. Our future work will focus on refining our data selection strategies to assemble a more diverse and robust pre-training corpus, as well as pursuing the directions discussed in Section \ref{sec:effec_rl} to further enhance reinforcement learning approaches for LLMs.

\textbf{Two-Stage Training}

\begin{itemize}[topsep=0pt]
    \item \textbf{High-quality}:
    We assess model capabilities on 8 mathematical benchmarks employing few-shot chain-of-thought prompting \cite{cot}. According to Table \ref{tab:corpora_comparison}, models trained with the DeepSeekMath~Corpus achieve superior outcomes. In Figure \ref{fig:corpora_comparisons}, it is evident that the DeepSeekMath-trained model consistently surpasses its Proof-Pile-2 counterpart at 50B tokens (equivalent to one complete pass over Proof-Pile-2), suggesting that DeepSeekMath delivers higher average data quality.
    \item \textbf{Multilingual}:
    Data within the DeepSeekMath~Corpus spans numerous languages, with English and Chinese comprising the largest portions. Table \ref{tab:corpora_comparison} illustrates that using the DeepSeekMath~Corpus for training bolsters mathematical reasoning for both English and Chinese. Conversely, traditional math datasets—largely English-focused—exhibit limited impact, and at times a detrimental effect, on Chinese mathematical reasoning.
    \item \textbf{Large-scale}:
    The DeepSeekMath~Corpus greatly exceeds existing mathematical datasets in terms of volume. Figure \ref{fig:corpora_comparisons} reveals that DeepSeek-LLM 1.3B demonstrates faster and more enduring gains in performance when trained with DeepSeekMath~Corpus. By contrast, the baseline corpora, which are much smaller, are quickly exhausted over repeated epochs, causing their corresponding model performance to stabilize early.
\end{itemize}

Subsequent to pre-training, we fine-tune DeepSeekMath-Base on mathematical instructions using approaches like chain-of-thought \citep{cot}, program-of-thought \citep{pot,pal}, and tool-integrated reasoning \citep{tora}. This procedure yields DeepSeekMath-Instruct 7B, which surpasses all other 7B-level models and rivals open-source instruction-tuned models in the 70B class.

To assess the impact of code training on mathematical reasoning proficiency, we examined both a two-stage and a single-stage training configuration:

Additionally, pretraining with code data benefits mathematical reasoning even in scenarios where tool use is unavailable. Within a two-stage pipeline, the initial code-focused training step already yields a noticeable uplift. Subsequent math training becomes more effective as a result, leading to peak performance. Conversely, mixing code and math tokens for one-stage training appears to diminish mathematical reasoning ability in tool-free contexts. A plausible explanation is that the representational capacity of DeepSeek-LLM 1.3B may be insufficient to accommodate the nuances of both programming and mathematical knowledge when learned simultaneously.

The models discussed are primarily designed for general applications, with most having undergone multiple alignment processes.

\begin{algorithm}[t]
  \small
  \caption{Iterative Group Relative Policy Optimization}
  \textbf{Input} initial policy model $\pi_{\theta_{\text{init}}}$; reward models $r_\varphi$; task prompts $\mathcal{D}$; 
  hyperparameters $\varepsilon$, $\beta$, $\mu$
  \begin{algorithmic}[1]
    \State Set policy model $\pi_\theta \leftarrow \pi_{\theta_{\text{init}}}$
    \For{iteration = 1, \dots, I}
       \State Assign reference model $\pi_{ref} \leftarrow \pi_{\theta}$
      \For{step = 1, \dots, M}
      \State Draw a batch $\mathcal{D}_b$ from $\mathcal{D}$
      \State Update previous policy model $\pi_{\theta_{old}} \leftarrow \pi_{\theta}$
      \State For each question $q \in \mathcal{D}_b$, generate $G$ outputs $\{o_i\}_{i=1}^G \sim \pi_{\theta_{old}} (\cdot \mid q) $
      \State Obtain rewards $\{r_i\}_{i=1}^{G}$ by evaluating each sampled output $o_i$ using $r_{\varphi}$
      \State Estimate $\hat{A}_{i,t}$ for the $t$-th token of $o_i$ utilizing group relative advantage calculation
      \For{GRPO iteration = 1, \dots, $\mu$}
        \State Refine policy model $\pi_{\theta}$ by maximizing the GRPO objective (Equation \ref{eq:GC-GRPO})
      \EndFor
    \EndFor \State Refine $r_\varphi$ through ongoing training with a replay buffer. 
    \EndFor 
  \end{algorithmic}
  \textbf{Output} $\pi_\theta$
  \label{alg:iter-grpo}
\end{algorithm}

\begin{itemize}[topsep=0pt]
    \item \textbf{Closed-source models} encompass: (1) members of the GPT series, in particular, GPT-4 \citep{gpt4} and GPT-4 Code Interpreter~\footnote{\url{https://openai.com/blog/chatgpt-plugins\#\#code-interpreter}}, (2) Gemini Ultra and Pro \citep{gemini}, (3) Inflection-2 \citep{inflection-2}, (4) Grok-1~\footnote{\url{https://x.ai/model-card}}, as well as recent offerings from Chinese organizations such as (5) Baichuan-3~\footnote{\url{https://www.baichuan-ai.com}}, and (6) the current GLM-4~\footnote{\url{https://open.bigmodel.cn/dev/api\#glm-4}} from the GLM line \citep{glm}.
\end{itemize}

\subsection{Lessons Learnt in Pre-Training}

Reinforcement learning (RL) has been demonstrated to further boost the mathematical problem-solving capacities of LLMs following the Supervised Fine-Tuning (SFT) stage \citep{wang2023math,wizardmath}. Here, we describe our RL method, Group Relative Policy Optimization (GRPO), which is designed to be both efficient and effective.

To assess the comparative quality of the DeepSeekMath~Corpus, we conduct pre-training studies alongside several recently released mathematical corpora:

\textbf{Exploration and Analysis of Reinforcement Learning}

We assemble an instruction-tuning dataset for mathematics, encompassing both English and Chinese tasks sourced from a range of mathematical domains and complexity levels. Each problem is accompanied by a solution articulated using chain-of-thought (CoT) \citep{cot}, program-of-thought (PoT) \citep{pot,pal}, or tool-integrated reasoning strategies \citep{tora}. The aggregate training corpus contains 776K instances.

    \item We provide insights gained through a series of math-focused training experiments. In particular, training models on code before proceeding with mathematical tasks enhances their ability to solve mathematical problems, with or without tool assistance. These observations contribute to addressing the persistent question: \textit{does code training enhance reasoning capacity?} Our findings suggest this is the case, especially concerning mathematical reasoning.

\paragraph{Mathematical Problem Solving with Step-by-Step Reasoning}

This work presents DeepSeekMath, a language model tailored to the mathematical domain that exhibits substantial advancements over open-source alternatives and demonstrates performance on par with GPT-4 across various academic benchmarks. Central to this progress is the construction of the DeepSeekMath Corpus—a comprehensive, high-caliber pre-training resource consisting of 120B mathematics-related tokens. The corpus is curated from the Common Crawl (CC) through the application of a fastText-based classifier \citep{joulin2016fasttext}. Initially, this classifier is trained using samples from OpenWebMath \citep{openwebmath} as positive cases, supplemented by a wide range of unrelated web content as negative cases. In the following phase, the classifier is deployed to identify further positive samples from the CC, which are then meticulously filtered and labeled by human annotators. These newly annotated samples are incorporated into the training set to enhance the classifier’s accuracy. Empirical evaluations reveal that the resulting large-scale dataset maintains a high standard of quality, as demonstrated by the DeepSeekMath-Base 7B model achieving 64.2\% accuracy on GSM8K \citep{gsm8k} and 36.2\% on the MATH competition dataset \citep{MATH}, surpassing Minerva 540B \citep{minerva}. Notably, the multilingual design of the DeepSeekMath Corpus contributes to performance gains on Chinese mathematics benchmarks \citep{wei2023cmath,agieval}. We regard our methodology for processing mathematical data as a foundational contribution to the field, with considerable opportunities remaining for further enhancement.

This work introduces DeepSeekMath, which establishes new state-of-the-art results among open-source systems on the competition-level MATH benchmark and approaches the accuracy of proprietary models. DeepSeekMath~is built on top of DeepSeek-Coder-v1.5 7B, followed by continual training over 500B tokens, including 120B mathematical tokens predominantly collected from Common Crawl. Through thorough ablation experiments, we demonstrate that web-sourced content offers considerable advantages as a source of mathematical data, whereas arXiv materials do not yield the anticipated benefits. We present Group Relative Policy Optimization (GRPO)—an alternative to Proximal Policy Optimization (PPO)—capable of markedly boosting mathematical reasoning performance while using less memory. Empirical evidence highlights GRPO’s effectiveness even as DeepSeekMath-Instruct 7B reaches high benchmark scores. Additionally, we propose a unified framework to analyze these techniques and summarize key areas where reinforcement learning for LLMs can be advanced.

Proximal Policy Optimization (PPO) \citep{schulman2017proximal} is a well-established actor-critic RL technique extensively adopted in the fine-tuning phase of LLMs \citep{ouyang2022training}. PPO guides the optimization of LLMs through maximization of the following surrogate objective function:

As detailed in Table \ref{tab:base_math_cot}, DeepSeekMath-Base 7B achieves top results across all eight evaluated benchmarks compared to other open-source base models, such as the widely adopted Mistral 7B \citep{mistral} and the recently introduced Llemma 34B \citep{llemma}, which received additional mathematical training using Proof-Pile-2 \citep{llemma}. Particularly on the competition-standard MATH dataset, DeepSeekMath-Base exceeds the best-performing open-source base models by more than 10\% in absolute terms, and even surpasses Minerva 540B \citep{minerva}—a proprietary base model that is 77 times larger, developed upon PaLM \citep{palm}, and further pretrained on mathematical corpora.

Here, both $\varepsilon$ and $\beta$ serve as tuning parameters. The quantity $\hat{A}_{i,t}$ stands for the advantage, but it is computed strictly from the relative rewards among the group of outputs for each question, a process discussed further in subsequent sections. The group-wise advantage calculation method employed by GRPO closely fits the inherently comparative nature of most reward models, which are often trained on datasets involving direct output comparisons for the same prompt. Furthermore, in contrast to modifying the reward with a KL term, GRPO penalizes deviation from the reference policy by directly adding a KL divergence regularizer between the new policy and the reference policy in the loss, thus keeping the computation of $\hat{A}_{i,t}$ straightforward. Unlike the KL regularization form in (\ref{eq:PPO-reward}), the KL divergence here is computed with the following unbiased estimator \citep{kl_approx}

Table \ref{tab:sft_rl_math} presents results under an evaluation scenario that prohibits the use of external tools. Here, DeepSeekMath-Instruct 7B achieves notably strong stepwise reasoning performance. In particular, on the advanced MATH dataset, this model outperforms every open-source competitor as well as most proprietary options (such as Inflection-2 and Gemini Pro) by a margin of at least 9 percentage points. This superiority holds even in comparison to models with considerably more parameters (for instance, Qwen 72B) or those leveraging mathematics-centric reinforcement learning strategies (like WizardMath-v1.1 7B). While DeepSeekMath-Instruct achieves results comparable to Chinese proprietary systems such as GLM-4 and Baichuan-3 on the MATH benchmark, it still trails behind GPT-4 and Gemini Ultra.

\paragraph{Algorithms} Algorithms utilize both data and reward signals to compute the gradient coefficient, which serves to adjust model parameters. According to Equation \ref{eq:objective}, contemporary methods largely place complete \textbf{TRUST} in the reward function's signal, using it to modulate the conditional probability of particular tokens. Nonetheless, the reliability of the reward signal cannot always be guaranteed, particularly for highly complex tasks. As an illustration, even rigorously curated datasets such as PRM800K \citep{lightman2023let}, annotated by skilled professionals, exhibit roughly 20\% annotation errors\footnote{\url{https://github.com/openai/prm800k/issues/12\#issuecomment-1728491852}}. Given these limitations, our focus turns to reinforcement learning algorithms engineered to withstand noisy or unreliable reward feedback. We contend that the adoption of \textbf{WEAK-TO-STRONG} \citep{burns2023weak} alignment techniques has the potential to fundamentally transform prevailing approaches to learning algorithms.

The constituent elements of these methodologies are outlined in Table \ref{tab:data_policy}. For a comprehensive exposition of the derivation procedure, see Appendix \ref{app:analysis-rl}.

\subsubsection{Outcome Supervision RL with GRPO} Given a question $q$, a set of $G$ candidate outputs $\{o_1, o_2, \cdots, o_G\}$ are generated from the old policy $\pi_{\theta_{old}}$. A reward model is applied to evaluate these outputs, producing a reward vector $\mathbf{r} = \{r_1, r_2, \cdots, r_G\}$. These scores are standardized by subtracting the group mean and dividing by the group standard deviation. Under outcome supervision, the normalized reward for each output $o_i$ is used as the advantage for all its tokens, i.e., $\hat{A}_{i, t} = \widetilde{r}_i = \frac{r_i- {\rm mean}(\mathbf{r})}{{\rm std}(\mathbf{r})}$ for every $t$, and policy optimization proceeds by maximizing the objective specified in equation (\ref{eq:GRPO-obj}).

\begin{itemize}[topsep=0pt]
    \item The potential influence of arXiv-sourced data on specialized mathematical tasks omitted from this study, such as informalizing theorem statements or proofs;
    \item The possible effects arising from arXiv data when jointly leveraged with other corpora;
    \item Whether scaling up to larger model architectures could reveal hitherto unseen benefits from arXiv-based pre-training.
\end{itemize}

\paragraph{Reflection on Gradient Coefficient} The procedure by which each algorithm transforms the reward signal into a gradient coefficient determines how the model parameters are updated. In our study, the reward function is constructed according to two schemes: `Rule’ and `Model’. Under `Rule’, the appropriateness of an answer is assessed directly by its correctness, while under `Model’, a dedicated reward model—trained using data labeled via the rule-based approach—assigns a score to each response. As delineated in Equations \ref{eq:GC-RFT} and \ref{eq:GC-GRPO}, GRPO diverges from Online RFT in that it tailors the gradient coefficient for each response based on the magnitude of the reward predicted by the reward model. This mechanism enables the algorithm to distinctly reinforce or penalize responses depending on the magnitude of their assessed value. Conversely, Online RFT does not incorporate such differentiation: it refrains from penalizing erroneous responses and applies uniform positive reinforcement to all responses judged correct, without modulating the strength of the update.

This section details the development of DeepSeekMath-Base 7B, a foundational model engineered for advanced mathematical reasoning. The model initialization leverages DeepSeek-Coder-Base-v1.5 7B \citep{deepseek-coder} and undergoes training on a corpus comprising 500B tokens. The data composition is as follows: 56\% sourced from the DeepSeekMath Corpus, 4\% from AlgebraicStack, 10\% from arXiv, 20\% from Github code, and an additional 10\% from Common Crawl natural language content in both English and Chinese. Training protocols align with those outlined in Section \ref{sec:quality-policy}, with adjustments made to employ a peak learning rate of 4.2e-4 and a token batch size of 10M.

\textbf{Math Pre-Training at Scale}

\begin{itemize}[topsep=0pt]
    \item We present Group Relative Policy Optimization (GRPO), a novel reinforcement learning approach designed for efficiency and effectiveness. GRPO eliminates the need for a critic model, relying instead on baseline estimations derived from group scores. This approach substantially decreases computational requirements during training compared to Proximal Policy Optimization (PPO).
    \item Our results show that employing GRPO substantially improves the performance of our instruction-tuned model, DeepSeekMath-Instruct, using only instruction-tuning data. Additionally, reinforcement learning with GRPO contributes to notable gains in out-of-domain generalization.
    \item We establish a unified framework that contextualizes diverse reinforcement learning techniques such as RFT, DPO, PPO, and GRPO. Through a series of thorough experiments—spanning online versus offline learning, outcome-based versus process-based supervision, and single-turn versus iterative reinforcement learning—we probe the core components critical to this paradigm.
    \item Building on this unified perspective, we delve into the underlying causes of reinforcement learning's effectiveness for large language models, and we outline several promising avenues for further enhancing LLM training via reinforcement learning.
\end{itemize}

\begin{itemize}
    \item \textbf{Supervised Fine-tuning (SFT)}: SFT further adapts a pretrained model by training on curated human-annotated SFT datasets.
    \item \textbf{Rejection Sampling Fine-tuning (RFT)}: RFT continues training the SFT model by utilizing only those outputs, sampled from the SFT model on the basis of its prompts, that are verified to be correct according to predetermined criteria.
    \item \textbf{Direct Preference Optimization (DPO)}: DPO extends SFT by training with additional responses sampled from the SFT model, employing a pairwise loss based on human preferences.
    \item \textbf{Online Rejection Sampling Fine-tuning (Online RFT)}: Unlike RFT, Online RFT commences with the SFT-initialized model, but generates and selects outputs dynamically from the continually updated policy model for further training.
    \item \textbf{PPO/GRPO}: These approaches begin from the SFT checkpoint and employ reinforcement learning, further adjusting the policy using live samples drawn from the evolving policy model itself.
\end{itemize}

\subsection{Training and Evaluating DeepSeekMath-Instruct 7B}

\subsection{Contributions}
This work presents advancements in large-scale math pre-training as well as an investigation and assessment of reinforcement learning methodologies.

\item \textbf{Open-source models} comprise both general-purpose models and those specifically tailored for enhanced mathematical abilities: (1) DeepSeek-LLM-Chat 67B \citep{deepseek-llm}, (2) Qwen 72B \citep{qwen}, (3) SeaLLM-v2 7B \citep{seallm}, and (4) ChatGLM3 6B \citep{chatglm3} represent the general models. Additionally, the mathematics-focused set includes (5) InternLM2-Math 20B~\footnote{\url{https://github.com/InternLM/InternLM-Math}}, derived from InternLM2 and specifically trained for mathematical tasks with further instruction tuning, (6) Math-Shepherd-Mistral 7B, which implements PPO training \citep{schulman2017proximal} on Mistral 7B \citep{mistral} utilizing a reward model based on process supervision, (7) WizardMath models \citep{wizardmath}, which advance mathematical problem solving in Mistral 7B and Llama-2 70B \citep{llama2} through evolve-instruct methods (AI-generated instructions in instruction tuning) and PPO training, predominantly using problems from GSM8K and MATH, (8) MetaMath 70B \citep{metamath}, which involves Llama-2 70B refined on enriched GSM8K and MATH datasets, (9) ToRA 34B \cite{tora}, a version of CodeLlama 34B adapted for tool-assisted mathematical reasoning, and (10) MAmmoTH 70B \citep{MathInstruct}, where Llama-2 70B is further instruction-tuned using MathInstruct.

\subsubsection{Why RL Works?} This study employs reinforcement learning on a subset of instruction tuning data and observes significant gains relative to the original instruction tuning model. To probe the mechanism behind RL’s effectiveness, we compare the Pass@K and Maj@K accuracies for both Instruct and RL models across two benchmarks. Figure \ref{fig:maj-pass} shows that RL substantially lifts Maj@K accuracy but does not impact Pass@K. These results suggest that the strength of RL lies in enhancing the robustness of the output distribution, \textbf{specifically by elevating the probability of selecting the correct answer from the TopK options rather than advancing the model’s core competencies.} Likewise, \citep{wang2023making} identified a \textbf{misalignment problem} in the reasoning tasks of SFT models and demonstrated that such models benefit from preference alignment strategies \citep{yuan2023rrhf,song2023preference,wang2023making} to improve reasoning ability.

    On the English benchmarks, DeepSeekMath-Base matches the performance of proprietary models such as Minerva 540B \citep{minerva} and exceeds the capabilities of all open-source base models, including Mistral 7B \citep{mistral} and Llemma-34B \citep{llemma}, regardless of their exposure to math-specific pre-training—and in many cases, by a considerable lead. Impressively, DeepSeekMath-Base attains even stronger results on the Chinese benchmarks, an outcome we attribute to our inclusion of diverse, high-quality non-English math data, contrasting with earlier methods \citep{minerva,llemma} that prioritized English-only corpora. Following mathematical instruction tuning and reinforcement learning, DeepSeekMath-Instruct and DeepSeekMath-RL achieve robust performance, with DeepSeekMath-RL crossing the 50\% accuracy threshold on the challenging MATH dataset—a first for open-source models.
\end{itemize}

\begin{itemize}[topsep=0pt]
    \item \textbf{Code Training for 400B Tokens $\rightarrow$ Math Training for 150B Tokens}: The DeepSeek-LLM 1.3B model is first exposed to 400B code tokens, followed by additional training on 150B math tokens.
    \item \textbf{General Training for 400B Tokens $\rightarrow$ Math Training for 150B Tokens}: As a baseline, we replace code tokens with general tokens (randomly selected from a large-scale general corpus constructed by DeepSeek-AI) during the initial 400B token training phase. This setup aims to evaluate whether pretraining with code tokens provides unique advantages in fostering mathematical reasoning capabilities compared to generic textual data.
\end{itemize}

\section{Supervised Fine-Tuning}

\begin{itemize}[topsep=0pt]
    \item We demonstrate that the openly available Common Crawl corpus holds substantial mathematical content. Leveraging a carefully crafted data selection framework, we curated the DeepSeekMath Corpus—an extensive, high-fidelity dataset comprising 120B tokens derived from web resources screened for mathematics. This corpus surpasses the math web dataset size used by Minerva \citep{minerva} by a factor of nearly 7 and is approximately 9 times larger than the recently introduced OpenWebMath \citep{openwebmath}.

\subsection{Training and Evaluating DeepSeekMath-Base 7B}

where $r_\varphi$ indicates the reward model, $\pi_{ref}$ is the reference model—commonly initialized from the SFT model—and $\beta$ scales the KL term.

\begin{itemize}[topsep=0pt]
    \item \textbf{MathPile} \citep{mathpile}: a dataset containing 8.9B tokens, predominantly (>85\%) composed of scientific arXiv articles, curated using specific cleaning and filtering heuristics;
    \item \textbf{ArXiv-RedPajama} \citep{redpajama}: a collection comprising 28.0B tokens from all available arXiv LaTeX source files, systematically purged of preambles, macros, comments, and bibliographies.
\end{itemize}

Accordingly, more extensive investigation is warranted and is left for future research efforts.

While there exists a prevalent but unconfirmed assertion that code training boosts reasoning capabilities, we seek to offer empirical evidence—specifically in the context of mathematics: code training enhances a model’s mathematical reasoning skills, with or without tool assistance.

Our reinforcement learning experiments use DeepSeekMath-Instruct 7B as the base model. For training, we draw from the SFT dataset's chain-of-thought-style items relevant to GSM8K and MATH, totaling about 144K samples. To more precisely study RL effects on benchmarks with limited or absent RL-stage supervision, we omit other SFT data. Reward model training data is constructed according to the methodology of \citep{wang2023math}. We initialize the reward model on DeepSeekMath-Base 7B using a learning rate of 2e-5. The policy model, under the GRPO paradigm, is updated with a learning rate of 1e-6 and a KL penalty set at 0.04. Each question prompts the generation of $64$ candidate responses. The model context and batch sizes are both set to 1024 tokens and instances, respectively. The policy receives a single optimization update for every exploration round. Evaluation of DeepSeekMath-RL 7B follows the same benchmarks as DeepSeekMath-Instruct 7B. For DeepSeekMath-RL 7B, GSM8K and MATH—when accompanied by chain-of-thought annotations—are considered in-domain, whereas the remaining benchmarks are treated as out-of-domain.

Because PPO’s value function is usually implemented as a model with similar capacity as the policy, this setup introduces considerable computational and memory overhead. Moreover, the value function in RL training acts as a baseline in the advantage computation to reduce variance; in the LLM domain, the reward model usually only supplies a score for the final token, making per-token value estimation more difficult. To overcome these drawbacks, we introduce Group Relative Policy Optimization (GRPO), illustrated in Figure \ref{fig:grpo}. GRPO forgoes an explicit value function: rather than using a learned baseline as in PPO, it instead computes the baseline as the mean reward over multiple sampled completions for each prompt. Specifically, for each question $q$, GRPO generates a collection of outputs $\{o_1, o_2, \cdots, o_G\}$ from $\pi_{\theta_{old}}$ and maximizes the following objective with respect to the policy parameters:

\subsubsection{How to Achieve More Effective RL?}
\label{sec:effec_rl}
Our findings indicate that RL is highly effective on mathematical reasoning benchmarks. Furthermore, we outline a unified conceptual framework that accommodates a variety of prominent training strategies. According to this framework, all such techniques can be viewed as either direct or approximate implementations of RL. Equation \ref{eq:objective} outlines three central elements: Data Source, Algorithm, and Reward Function. We propose several promising research avenues for each component.

DeepSeekMath-Base is built upon DeepSeek-Coder-Base-v1.5 7B \citep{deepseek-coder}, as our findings indicate that initializing with a code-centric model leads to superior outcomes compared to beginning with a standard LLM. Notably, we find that math-centric pre-training also brings notable gains on MMLU \citep{mmlu} and BBH benchmarks \citep{bbh}, suggesting that improvements extend beyond mathematical proficiency to broader reasoning abilities.

\subsubsection{Process Supervision RL with GRPO} While outcome supervision yields a reward only at the end of each output, this may limit policy guidance in intricate mathematical reasoning. Motivated by \cite{wang2023math}, process supervision is introduced to deliver feedback at the end of each reasoning step. Specifically, for a given question $q$ and corresponding $G$ sampled outputs $\{o_1, o_2, \cdots, o_G\}$, a process reward model assesses each step of every output, providing rewards structured as $\mathbf{R} = \{ \{r_1^{index(1)},\cdots,r_1^{index(K_1)}\}, \cdots,  \{r_G^{index(1)},\cdots,r_G^{index(K_G)}\} \}$. Here, $index(j)$ marks the terminal token of the $j$-th step and $K_i$ counts the total steps in the $i$-th output. All rewards are normalized by their overall mean and standard deviation: $\widetilde{r}_i^{index(j)} = \frac{r_i^{index(j)} - {\rm mean(\mathbf{R})}}{{\rm std(\mathbf{R})}}$.
In process supervision, the advantage assigned to each token corresponds to the cumulative normalized reward over future steps, formally $\hat{A}_{i, t} = \sum_{index(j) \ge t} \widetilde{r}_i^{index(j)}$,
and policy parameters are updated via maximization of the objective in equation (\ref{eq:GRPO-obj}).

\begin{equation}
\footnotesize
\begin{split}
    \mathcal{J}_{GRPO}(\theta) &= \mathbb{E}{[q \sim P(Q), \{o_i\}_{i=1}^G \sim \pi_{\theta_{old}}(O|q)]}  \\
    & \frac{1}{G}\sum_{i=1}^G\frac{1}{|o_i|} \sum_{t=1}^{|o_i|} \left\{ \min \left[ \frac{\pi_\theta(o_{i,t} | q, o_{i,<t})}{\pi_{\theta_{old}}(o_{i,t} | q, o_{i,<t})} \hat{A}_{i,t}, \text{clip} \left( \frac{\pi_\theta(o_{i,t} | q, o_{i,<t})}{\pi_{\theta_{old}}(o_{i,t} | q, o_{i,<t})}, 1 - \varepsilon, 1 + \varepsilon \right)  \hat{A}_{i,t} \right] - \beta \mathbb{D}_{KL}\left[\pi_{\theta} || \pi_{ref}\right]\right\} ,
\end{split}
\label{eq:GRPO-obj}
\end{equation}

\section{Conclusion, Limitation, and Future Work} 

\begin{equation}
\footnotesize
    \mathcal{J}_{PPO}(\theta) = \mathbb{E}{[q \sim P(Q), o \sim \pi_{\theta_{old}}(O|q)]} \frac{1}{|o|} \sum_{t=1}^{|o|} \min \left[ \frac{\pi_\theta(o_{t} | q, o_{<t})}{\pi_{\theta_{old}}(o_{t} | q, o_{<t})} A_{t}, \text{clip} \left( \frac{\pi_\theta(o_{t} | q, o_{<t})}{\pi_{\theta_{old}}(o_{t} | q, o_{<t})}, 1 - \varepsilon, 1 + \varepsilon \right)  A_{t} \right] ,
\end{equation}

Automating the formalization of mathematical proofs not only promotes rigor and dependability but also improves productivity—an area that has drawn growing interest. DeepSeekMath-Base 7B is evaluated on the informal-to-formal proof conversion task described in \citep{dsp_proof}, which requires generating a formal proof from an informal statement, its formal version, and an accompanying informal proof. Our experiments on miniF2F \citep{minif2f}—a challenging formal mathematics benchmark at the Olympiad level—involve using few-shot prompting to produce Isabelle-formalized proofs. Adopting the approach from \cite{dsp_proof}, the models first generate proof sketches, and then the Sledgehammer tool \citep{sledgehammer} is applied to fill in detailed reasoning steps. Table \ref{tab:base_math_pot_proof} indicates that DeepSeekMath-Base 7B demonstrates strong competency in the autoformalization of proofs.

\subsubsection{Code Training Benefits Mathematical Reasoning}

Nonetheless, it is important to recognize several caveats regarding this finding. Notably, we have yet to address:

\section{Discussion}
This section details the insights derived from our pre-training and reinforcement learning experiments.

Table \ref{tab:sft_rl_math} presents the comparative results for open- and closed-source models employing both chain-of-thought and tool-augmented reasoning approaches on English and Chinese assessment sets. Our observations reveal the following: (1) Leveraging chain-of-thought reasoning, DeepSeekMath-RL 7B achieves accuracy rates of 88.2\% on GSM8K and 51.7\% on MATH, outperforming all open-source models within the 7B to 70B parameter scale, as well as exceeding the capabilities of most closed-source counterparts. (2) Significantly, DeepSeekMath-RL 7B's training regime consists exclusively of chain-of-thought-format instruction-tuning data drawn from GSM8K and MATH, using DeepSeekMath-Instruct 7B as its initialization point. Even with this limited training data, the model consistently surpasses DeepSeekMath-Instruct 7B across every evaluation dimension, underscoring the value added by reinforcement learning.

The model's proficiency in mathematical tasks is assessed in both tool-free and tool-augmented settings, using four quantitative reasoning evaluation sets across English and Chinese. We perform a comparative analysis with contemporaneous state-of-the-art models:

\subsection{Validating the Quality of the DeepSeekMath~Corpus}

\begin{itemize}[topsep=0pt]
    \item \textbf{English mathematical datasets}: We enhance GSM8K and MATH datasets by providing tool-integrated solutions, and incorporate selected items from MathInstruct \citep{MathInstruct} in addition to examples from the Lila-OOD \citep{lila} training split, with solutions presented in either CoT or PoT frameworks. Our English dataset spans a variety of mathematical areas, such as algebra, probability, number theory, calculus, and geometry.
    \item \textbf{Chinese mathematical datasets}: Our collection of Chinese-language problems covers K-12 curricula and 76 distinct subfields, for example, linear equations; for these, solutions are annotated utilizing both CoT and tool-integrated methodologies.
\end{itemize}

\subsubsection{Evaluation Results}

The advent of large language models (LLM) has significantly transformed mathematical reasoning within artificial intelligence, driving notable improvements in both quantitative reasoning \citep{MATH} and geometric reasoning benchmarks \citep{trinh2024solving}. These models have also demonstrated substantial utility in supporting humans with intricate mathematical problem-solving tasks \citep{tao}. Nevertheless, state-of-the-art systems such as GPT-4 \citep{gpt4} and Gemini-Ultra \citep{gemini} are not accessible to the public, and open-source alternatives presently lag markedly in terms of performance.

\paragraph{Formal Mathematics}

\paragraph{Mathematical Problem Solving with Tool Use} For program-assisted mathematical reasoning, we assess performance on GSM8K and MATH using few-shot program-of-thought prompting strategies \citep{pot,pal}. In this setup, models address each problem by constructing a Python program, making use of modules like \textit{math} and \textit{sympy} for advanced calculations. The computed result of the generated program is then used as the final answer. According to the results in Table \ref{tab:base_math_pot_proof}, DeepSeekMath-Base 7B sets a new state-of-the-art, outperforming the previous leading Llemma 34B.

\paragraph{Reward Function} The reward function constitutes the main training signal source. In reinforcement learning settings, it is typically implemented as a neural reward model. Our perspective is that three principal research avenues exist for reward models: 1) \textbf{Enhancing the generalization capabilities of the reward model.} For reinforcement learning to transcend mere stabilization of large language model distributions and to facilitate genuine capability improvement, reward models must be able to generalize to novel or out-of-distribution prompts and sophisticated outputs; 2) \textbf{Capturing the uncertainty associated with the reward model.} Modeling this uncertainty could bridge the gap between underperforming reward functions and weak-to-strong learning strategies; 3) \textbf{Constructing high-quality process reward models efficiently} to deliver granular, process-level feedback for reasoning tasks \citep{lightman2023let,wang2023math}.

\item \textbf{Formal Mathematics}: The performance of DeepSeekMath-Base is assessed on the informal-to-formal theorem proving task as outlined in \citep{dsp_proof}, utilizing miniF2F \citep{minif2f} as the benchmark and Isabelle \citep{isabelle} as the designated proof assistant. DeepSeekMath-Base achieves robust few-shot autoformalization outcomes.
\item \textbf{Natural Language Understanding, Reasoning, and Code}: To thoroughly gauge the general abilities of the models in understanding, reasoning, and programming, DeepSeekMath-Base is benchmarked on Massive Multitask Language Understanding (MMLU) \citep{mmlu}, which includes 57 multiple-choice tasks from varied fields; BIG-Bench Hard (BBH) \citep{bbh}, comprising 23 tasks that typically necessitate complex multi-step reasoning; as well as HumanEval \citep{codex} and MBPP \citep{mbpp}, both standard measures for assessing code generation models. Pre-training on mathematical data enhances both linguistic comprehension and reasoning efficacy.

\textbf{The DeepSeekMath~Corpus stands out for its exceptional quality, extensive multilingual representation, and unmatched scale.}

\subsection{SFT Data Curation}

Additionally, we present Group Relative Policy Optimization (GRPO), a novel reinforcement learning (RL) variant inspired by Proximal Policy Optimization (PPO) \citep{schulman2017proximal}. Unlike traditional approaches, GRPO omits the critic model and estimates the baseline directly from grouped score distributions, leading to substantial resource savings during training. Utilizing only a portion of English instruction data, GRPO achieves remarkable gains over the already strong DeepSeekMath-Instruct, both within domain (GSM8K: 82.9\% $\rightarrow$ 88.2\%, MATH: 46.8\% $\rightarrow$ 51.7\%) and in out-of-domain mathematics challenges (e.g., CMATH: 84.6\% $\rightarrow$ 88.8\%) during RL optimization. 

Although arXiv papers are frequently incorporated into mathematical pre-training datasets \citep{minerva,gpt-f,llemma,mathpile}, there has been little comprehensive evaluation of their specific influence on mathematical reasoning capabilities. Our findings suggest, perhaps surprisingly, that arXiv papers contribute minimally to advancements in mathematical reasoning. To assess this, we utilized DeepSeek-LLM 1.3B and DeepSeek-Coder-Base-v1.5 7B \citep{deepseek-coder} models and subjected them to pre-training on arXiv-based corpora, which were constructed through distinct processing workflows:

Referencing Figure \ref{fig:rl-analysis}, our empirical observations demonstrate that Online RFT achieves markedly superior performance relative to RFT across both evaluated benchmarks. In particular, during initial training phases, Online RFT exhibits results on par with RFT, yet it establishes a clear advantage as training progresses, illustrating the effectiveness of leveraging online sampling. This phenomenon can be rationalized by noting that, at early stages, the policy (actor) and the SFT model are largely similar, resulting in minimal distinction in the data generated. As training advances, however, the divergence between the actor’s policy and SFT increases, with online sampling yielding more informative and diverse examples, thus conferring substantial benefits.

We describe the instruction-tuning process of DeepSeekMath-Instruct 7B, which is initialized from DeepSeekMath-Base. During training, samples are randomly concatenated until the combined sequence reaches up to 4K tokens in length. The training regimen consists of 500 steps, with a batch size of 256, and utilizes a fixed learning rate of 5e-5.

\begin{itemize}[topsep=0pt]
    \item \textbf{English and Chinese Mathematical Reasoning}: Our evaluation covers an extensive set of benchmarks in both English and Chinese, addressing mathematical reasoning tasks from elementary through collegiate difficulty. The English suite includes GSM8K \citep{gsm8k}, MATH \citep{MATH}, SAT \citep{llemma}, OCW Courses \citep{minerva}, and MMLU-STEM \citep{mmlu}, while Chinese datasets encompass MGSM-zh \citep{mgsm}, CMATH \citep{wei2023cmath}, Gaokao-MathCloze \citep{agieval}, and Gaokao-MathQA \citep{agieval}. We assess models’ abilities to provide fully self-contained text-based solutions (without external tools), as well as their capabilities when solving via Python execution.

\subsection{Summary of Evaluations and Metrics}

\subsection{Data Collection and Decontamination} This section describes how the DeepSeekMath Corpus was curated from Common Crawl. Figure \ref{fig:data_collect} illustrates the iterative procedure developed to collect a large-scale mathematical dataset from Common Crawl, initiating from a compact, high-quality seed dataset focused on mathematics. The methodology can readily be transferred to domains beyond mathematics, such as coding.

\subsection{Group Relative Policy Optimization}

To initiate, we select OpenWebMath \citep{openwebmath}, a vetted collection of mathematical texts from the web, to serve as the foundational seed corpus. A fastText model \citep{joulin2016fasttext} is then trained to identify additional web pages with characteristics similar to those in OpenWebMath. For this, 500,000 examples are drawn at random from the seed corpus as positive samples, while another 500,000 pages from Common Crawl are utilized as negatives. Training is conducted with an open-source implementation\footnote{\url{https://fasttext.cc}} using a 256-dimensional vector space, a learning rate of 0.1, a maximum n-gram length of 3, a minimum word frequency of 3, and three training epochs. The initial Common Crawl dataset is reduced in size by applying URL-level deduplication and near-duplicate removal, yielding a set of 40 billion HTML web documents. Using the fastText classifier, mathematical pages are identified from this deduplicated dataset. To eliminate subpar mathematical content, retrieved pages are ranked by their fastText-assigned score, with only the highest-ranked examples retained. The effect of different dataset scales is evaluated by pre-training on subsets containing the top 40B, 80B, 120B, and 160B tokens, with the initial iteration focusing on the top 40B tokens.

To gauge the model's proficiency in mathematics, we implement a multifaceted evaluation addressing its ability to independently generate complete mathematical solutions, solve mathematical tasks utilizing various tools, and engage in formal theorem proving. Furthermore, we assess broader capabilities beyond mathematics, including natural language comprehension, reasoning faculties, and programming proficiency.

\begin{equation}
    r_{t} = r_\varphi(q, o_{\le t}) - \beta \log\frac{\pi_{\theta}(o_{t}|q, o_{<t})}{\pi_{ref}(o_{t}|q, o_{<t})},
\label{eq:PPO-reward}
\end{equation}

Incorporating code training enhances the model's ability to perform program-aided mathematical problem solving, irrespective of whether training is staged or conducted in a single phase. Table \ref{tab:code-math} reveals that, in a two-phase approach, code pretraining substantially improves performance on Python-based GSM8K and MATH tasks, with math-focused fine-tuning providing additional gains. Furthermore, when both code and math data are blended in a one-stage setting, the model not only avoids the catastrophic forgetting observed in standard two-stage pipelines, but also exhibits synergistic improvements in both code (Table \ref{tab:code-math-general-eval}) and program-assisted mathematical reasoning (Table \ref{tab:code-math}) tasks.

\subsubsection{Training Setting}
\label{sec:quality-policy}
We perform specialized math-oriented training on a pre-existing language model comprising 1.3B parameters, following the architecture employed by the DeepSeek LLMs \citep{deepseek-llm}, referred to here as DeepSeek-LLM 1.3B. Separate training is performed on each mathematical dataset for a total of 150B tokens. All experiments utilize the streamlined HAI-LLM \citep{haillm} training framework. Consistent with DeepSeek LLMs' training approach, optimization is carried out with AdamW \citep{adamW} using $\beta_1=0.9$, $\beta_2=0.95$, and $\mathrm{weight\_decay}=0.1$. The learning rate schedule features several stages: it peaks following 2,000 warmup steps, falls to 31.6\% of the maximum after 80\% of the training duration, and further declines to 10.0\% after 90\% completion. The highest learning rate is capped at 5.3e-4. Training employs a batch size of 4M tokens and a context length of 4K.

\subsubsection{Iterative RL with GRPO} During the course of reinforcement learning, the previously trained reward model may become inadequate for guiding the continuously evolving policy model. To address this, we investigate an iterative approach to RL using GRPO. As depicted in Algorithm \ref{alg:iter-grpo}, the iterative GRPO process entails periodically creating new datasets for reward model training based on trajectories sampled from the latest policy model. The reward model is updated incrementally, utilizing a replay buffer that retains 10\% of prior examples. Subsequently, the current policy model is adopted as the new reference model, and training proceeds with the refreshed reward model supervising the policy.

    \item Our DeepSeekMath-Base 7B base model, after pre-training, matches the performance of Minerva 540B \citep{minerva}, suggesting that scale alone does not account for mathematical reasoning proficiency. This result highlights that even models with a modest parameter count can reach strong performance when exposed to high-quality, domain-relevant data.

\paragraph{Natural Language Understanding, Reasoning, and Code}

\section{Math Pre-Training}

\paragraph{Data Source} The data source serves as the foundation for all training strategies. For RL, this specifically refers to unlabeled problems paired with solutions sampled from the current policy. In our experiments, we rely solely on questions derived from the instruction tuning dataset and employ basic nucleus sampling to obtain outputs. We hypothesize that this may explain why our RL pipeline selectively improves Maj@K. Moving forward, we plan to extend our RL framework to include out-of-distribution question prompts, alongside \textbf{more sophisticated sampling (decoding) approaches}, such as those employing tree-search methodologies \citep{yao2023tree}. Additionally, \textbf{efficient inference strategies} \citep{xia-etal-2023-speculative,leviathan2023fast,kwon2023efficient,xia2024unlocking}, which are crucial for boosting the exploration capabilities of policy models, represent another important direction.

Moreover, we introduce a cohesive framework for analyzing multiple methodologies such as Rejection Sampling Fine-Tuning (RFT) \citep{yuan2023scaling}, Direct Preference Optimization (DPO) \citep{dpo}, PPO, and GRPO. Through this unified perspective, we observe that these approaches can all be interpreted as variations of direct or streamlined RL algorithms. We systematically examine critical aspects of this paradigm, including comparisons between online and offline training, supervision at the output versus process level, as well as single-turn and iterative RL setups. Finally, we discuss the mechanisms through which our RL approach enhances instruction-tuned model performance and outline prospective avenues for further improving RL within this integrated framework.

To prevent contamination of evaluation benchmarks, we implement the filtering strategy outlined in \cite{deepseek-coder}. Web pages containing either questions or answers sourced from prominent English math benchmarks such as GSM8K \citep{gsm8k}, MATH \citep{MATH}, as well as Chinese benchmarks like CMATH \citep{wei2023cmath} and AGIEval \citep{agieval}, are removed. Specifically, any segment in the text corpus that contains a 10-gram sequence exactly matching any substring from the evaluation benchmarks is excluded from the training data. For benchmark texts with fewer than 10 grams but at least 3 grams, an exact match approach is utilized to eliminate web pages with potential overlap.

\begin{equation}
    \nabla_{\theta}\mathcal{J}_{\textcolor{red}{\mathcal{A}}}(\theta) = \mathbb{E}[\underbrace{(q,o) \sim \textcolor{red}{\mathcal{D}}}_{Data \ Source}]\left( \frac{1}{|o|} \sum_{t=1}^{|o|}  \underbrace{GC_{{\mathcal{A}}}(q, o, t, \textcolor{red}{\pi_{{rf}}})}_{Gradient \ Coefficient}  \nabla_{\theta}\log \pi_{\theta}(o_t | q, o_{<t})\right).
\label{eq:objective}
\end{equation}

\textbf{One-Stage Training}

We conducted isolated training sessions, involving 150B tokens for DeepSeek-LLM 1.3B and 40B tokens for DeepSeek-Coder-Base-v1.5 7B, for each respective arXiv dataset. Our results show that pre-training solely on arXiv data does not yield measurable improvements in mathematical reasoning performance; in certain cases, the models even underperform on a broad range of mathematical evaluation tasks featured in this work. Assessed benchmarks span various difficulties and include the GSM8K and MATH datasets for quantitative reasoning (see Table \ref{tab:arxiv-cot}), the MMLU-STEM multiple-choice suite (Table \ref{tab:arxiv-cot}), and the miniF2F set for formal mathematics (Table \ref{tab:arxiv_atp}).

We analyze DeepSeekMath-Base’s aptitude for mathematical problem-solving using few-shot chain-of-thought prompting \citep{cot}. The evaluation spans eight benchmarks in both English and Chinese, covering quantitative reasoning datasets such as GSM8K \citep{gsm8k}, MATH \citep{MATH}, and CMATH \citep{wei2023cmath}, as well as multiple-choice datasets like MMLU-STEM \citep{mmlu} and Gaokao-MathQA \citep{agieval}. These benchmarks collectively address a wide spectrum of mathematical topics, from elementary to collegiate level, emphasizing diverse aspects of mathematical reasoning.

\subsubsection{From PPO to GRPO}

For evaluating language understanding, we consider MMLU \citep{mmlu}; for reasoning, we assess BBH \citep{bbh}; and for coding proficiency, we include HumanEval \citep{codex} and MBPP \citep{mbpp}. Results in Table \ref{tab:understanding_reasoning_code} show that DeepSeekMath-Base 7B substantially improves upon the performance of its predecessor, DeepSeek-Coder-Base-v1.5 \citep{deepseek-coder}, on both MMLU and BBH, underscoring the effectiveness of mathematics-focused training for language and reasoning capabilities. Moreover, continual exposure to code tokens enables DeepSeekMath-Base 7B to maintain its predecessor’s high performance on the two programming benchmarks. Collectively, DeepSeekMath-Base 7B demonstrates significant superiority over the generalist Mistral 7B \citep{mistral} across the reasoning and coding assessments.

We begin by summarizing our pre-training experiences. Unless otherwise indicated, all experiments employ the training protocol specified in Section \ref{sec:quality-policy}. For clarity, references to the DeepSeekMath Corpus within this discussion pertain to the 89B-token version acquired during the second iteration of dataset construction.

\paragraph{Reflection on Data Source} We categorize the sources of data into two primary types: online sampling and offline sampling. Online sampling indicates that training examples are generated via the real-time policy model’s ongoing exploration, whereas offline sampling involves utilizing data gathered from the initial SFT model’s outputs. The RFT and DPO algorithms adhere to offline data acquisition, whereas Online RFT and GRPO utilize data collected from the current training policy—characteristic of the online paradigm.

When the evaluation allows the combination of natural language and code-based reasoning tools, DeepSeekMath-Instruct 7B attains an accuracy close to 60\% on MATH, setting a new high among all known open-source models. Across other benchmarks, its performance matches that of DeepSeek-LLM-Chat 67B—the previous open-source state-of-the-art, which has a model size ten times larger.

As illustrated in Figure \ref{fig:rl-analysis}, GRPO outperforms online RFT, underscoring the effectiveness of modifying the coefficients for positive and negative gradients. Moreover, the results show that GRPO+PS achieves better outcomes than GRPO+OS, demonstrating the advantage of employing more granular, step-sensitive gradient coefficients. Additionally, we investigate the use of iterative RL by conducting two rounds in our experiments. The results, depicted in Figure \ref{fig:iter-rl}, reveal that iterative RL leads to a marked increase in performance, particularly notable in the initial iteration.

    \item Despite the prevalence of arXiv datasets in math-oriented research, our experiments reveal that incorporating arXiv papers yields negligible benefit across the mathematical benchmarks evaluated in this study.
\end{itemize}

Here, $\pi_{\theta}$ denotes the current policy model, while $\pi_{\theta_{old}}$ refers to the previous policy version; the variables $q$ and $o$ represent questions and generated outputs, which are sampled from both the question dataset and the older policy $\pi_{\theta_{old}}$. The parameter $\varepsilon$ is a clipping hyper-parameter from PPO, serving to stabilize optimization. The term $A_t$ designates the advantage, estimated using Generalized Advantage Estimation (GAE) \citep{gae} based on rewards $\{r_{\ge t}\}$ and a learned value function $V_{\psi}$. Consequently, PPO requires simultaneous training of a value function and the policy itself. To prevent the reward model from being over-optimized, the common practice is to introduce a per-token KL penalty between the policy and a reference model as an additive term to the reward at each token \citep{ouyang2022training}, formulated as:

Following the initial round of data acquisition, a significant number of mathematical web pages remain uncollected. This is primarily due to the limited diversity of positive examples used during the training of the fastText model. To address this issue and enhance the representativeness of the seed corpus, we expand our collection by sourcing additional mathematical web domains. Specifically, the entire Common Crawl dataset is partitioned into distinct domains, where each domain is characterized by a shared base URL among its web pages. For each domain, we determine the proportion of pages that were retrieved in the initial collection round. Domains in which more than 10\% of web pages are retrieved are designated as math-related (for instance, \url{mathoverflow.net}).
Subsequent to this step, we manually review and annotate those URLs within the identified domains that pertain to mathematical content (e.g., \url{mathoverflow.net/questions}). Any yet-uncollected web pages associated with these annotated URLs are appended to the seed corpus. Through this process, we accumulate a greater number of positive samples, allowing us to retrain the fastText model to more effectively capture mathematical content during further collection. After completing four cycles of data collection, the result is a set of 35.5M mathematical web pages comprising a total of 120B tokens. Notably, during the fourth cycle, it is observed that almost 98\% of these data had already been gathered by the conclusion of the third cycle, prompting the cessation of further collection.

\begin{itemize}[topsep=0pt]
    \item \textbf{Math Training for 150B Tokens}: Here, DeepSeek-LLM 1.3B is trained exclusively on 150B math tokens from scratch.
    \item \textbf{Training on a mixture of 400B Code Tokens and 150B Math Tokens}: Sequential math training after code pretraining can impair code-related performance. Thus, we examine whether a simultaneous exposure to both code and math tokens during a single-stage training regime not only bolsters mathematical reasoning but also counteracts catastrophic forgetting.
\end{itemize}

\subsubsection{Limited Effectiveness of ArXiv Papers in Enhancing Mathematical Reasoning}

\paragraph{Results}
Tables \ref{tab:code-math} and \ref{tab:code-math-general-eval} summarize model outcomes under each of the aforementioned training configurations.

\subsection{Training and Evaluation of DeepSeekMath-RL}

This formulation consists of three principal elements: 1) \textit{Data Source $\mathcal{D}$}, which supplies the input data for training; 2) \textit{Reward Function $\pi_{{rf}}$}, serving as the mechanism that assigns reward signals; and 3) \textit{Algorithm $\mathcal{A}$}, which utilizes both the training data and reward signals to calculate the gradient coefficient $GC$, determining the intensity of positive or negative updates for each example. Within this general framework, we compare several widely adopted methods:

\subsection{Reinforcement Learning: Emerging Insights}
\subsubsection{Towards a Unified Framework} This section proposes a common framework for comparing various training approaches—including SFT, RFT, DPO, PPO, and GRPO—and presents empirical findings on critical factors shaping this framework. Typically, the parameter gradient $\theta$ for these training algorithms can be formalized as:

\section{Reinforcement Learning}

\section{Introduction}

\begin{itemize}[topsep=0pt]
    \item \textbf{MathPile} \citep{mathpile}: This multi-source dataset encompasses 8.9B tokens drawn from resources such as textbooks, Wikipedia, ProofWiki, CommonCrawl, StackExchange, and arXiv, with arXiv contributions accounting for more than 85\% of the corpus;
    \item \textbf{OpenWebMath} \citep{openwebmath}: Composed of CommonCrawl records specifically filtered for mathematical relevance, comprising 13.6B tokens in total;
    \item \textbf{Proof-Pile-2} \citep{llemma}: Constructed from OpenWebMath, AlgebraicStack (offering 10.3B tokens of mathematical code), and arXiv documents (providing 28.0B tokens). For experimentation on Proof-Pile-2, we adhere to the arXiv:Web:Code sampling ratio of 2:4:1, following \cite{llemma}.
\end{itemize}